\documentclass[tikz,border=6pt]{standalone}
\usepackage{peabra}
\begin{document}
\begin{tikzpicture}[x=1cm,y=1cm,
  csp/.style={doboz=34mm, font=\footnotesize\bfseries, minimum height=12mm},
  tev/.style={-{Stealth[length=2.5mm]}, line width=1pt, draw=peSotet},
  vv/.style={font=\scriptsize\bfseries, text=peZold, align=center, fill=white, inner sep=1.5pt}]
  \node[sotetdoboz=34mm, font=\footnotesize\bfseries, minimum height=12mm] (p) at (0,3.2) {a valós rendszer (problémaentitás)};
  \node[csp] (c) at (-4.4,-1.4) {fogalmi modell};
  \node[csp] (s) at (4.4,-1.4) {számítógépes modell};
  \draw[tev] (p.west) to[bend right=25] node[font=\scriptsize, left, align=right, pos=0.4] {elemzés és\\modellalkotás} (c.north);
  \draw[tev] (c.east) -- node[font=\scriptsize, below, align=center] {programozás és\\megvalósítás} (s.west);
  \draw[tev] (s.north) to[bend right=25] node[font=\scriptsize, right, align=left, pos=0.6] {kísérletezés\\a modellel} (p.east);
  \node[vv] at (-1.75,1.35) {fogalmi modell\\validitása};
  \node[vv] at (0,-0.85) {számítógépes modell\\verifikációja};
  \node[vv] at (1.75,1.35) {működési\\validitás};
  \node[kiemeltdoboz=28mm, font=\scriptsize] at (0,0.35) {adatvaliditás};
\end{tikzpicture}
\end{document}
